\section{Discussão}
\label{sec:discussao}

Os resultados mostram convergência em torno de uma compreensão ampliada da manutenção predial. Desempenho técnico, governança, ambiente, ciclo de vida, custos e condições sociais de uso aparecem de forma combinada, o que aproxima a manutenção de uma função de gestão de ativos e de facilidades, e não apenas de resposta corretiva. Essa convergência é compatível com a literatura de gestão sustentável de facilidades, que distribui responsabilidades entre os níveis estratégico, tático e operacional e relaciona a operação dos edifícios aos objetivos institucionais \parencite{nielsen_sustentabilidadefm_2016,opoku_futurofm_2022,okoro_sustainablefm_2023}. No corpus, entretanto, a frequência elevada das dimensões técnica-operacional e institucional decorre de categorias documentais amplas. Ela indica recorrência temática, mas não mede a qualidade ou a intensidade com que cada dimensão foi operacionalizada.

A distribuição temporal evidencia expansão recente da produção, com concentração de 84,6\% dos registros entre 2019 e 2026 e maior frequência em 2025. Esse padrão é contemporâneo à ampliação do uso de BIM, \textit{digital twins}, sensores, modelos de informação e métodos de otimização na operação de edifícios. A associação temporal, contudo, deve ser interpretada de forma descritiva, pois a revisão não avaliou causalidade entre difusão tecnológica e volume de publicações. A literatura examinada em texto completo ilustra diferentes estágios de maturidade: há revisões que ainda identificam distância entre BIM e \textit{digital twins} plenamente integrados, ao mesmo tempo que surgem modelos de informação de ativos voltados à manutenção preditiva \parencite{deng_bimdigitaltwins_2021,gispert_ontologiaativos_2025,goretti_bimoem_2023}.

Também há divergência entre a ampla presença de abordagens gerais e a baixa frequência de métodos formalmente estruturados. A categoria \textit{framework} aparece em quase todo o núcleo, enquanto AHP, TOPSIS, ANP, Delphi e MCDM ocorrem em subconjuntos reduzidos. Essa diferença indica que propor uma estrutura conceitual ou mencionar apoio à decisão não equivale a aplicar um procedimento multicritério com critérios definidos, normalização, ponderação, agregação, análise de sensibilidade e validação. Nos casos em que a aplicação está documentalmente sustentada, os métodos assumem funções distintas: avaliação de sustentabilidade de edifícios educacionais, construção de referenciais operacionais de \textit{campus} e alocação de prestadores de manutenção de edifícios públicos \parencite{alfalah_frameworkeducacional_2022,naji_campusdelphi_2023,masmoudi_ahptopsis_2026}.

A mesma cautela se aplica aos critérios de priorização. Desempenho operacional, informação e dados, custo, vida útil, energia, risco e condição física formam um núcleo recorrente, mas a presença desses termos não demonstra que tenham sido convertidos em indicadores mensuráveis ou efetivamente usados em uma decisão. Parte dos estudos se mantém no plano conceitual, enquanto estudos empíricos e quantitativos operacionalizam subconjuntos específicos conforme o objeto analisado. Modelos de custo para edifícios educacionais, planejamento conjunto de manutenção e desempenho energético e simulações estocásticas de estratégias de intervenção exemplificam níveis mais explícitos de operacionalização \parencite{kim_estimativacustoseducacional_2018,nageli_ciclovida_2019,ferreira_manutencaoceramica_2020}. Assim, a matriz proposta deve preservar a distinção entre critério identificado, variável mensurada e critério efetivamente aplicado.

Os contextos de aplicação revelam outra assimetria. Edificações genéricas e portfólios prediais predominam, enquanto universidades, \textit{campus} e edifícios públicos representam parcelas menores. A categoria edificação genérica funciona como categoria-guarda-chuva e não caracteriza uma tipologia homogênea. Já os estudos voltados a instituições de ensino destacam fatores que dificilmente podem ser tratados apenas como desempenho físico: diversidade de usuários, continuidade das atividades acadêmicas, restrições orçamentárias, governança, experiência dos ocupantes e articulação entre operação e objetivos educacionais \parencite{awuzie_fmuniversidades_2017,alfalah_frameworkeducacional_2022,naji_campusdelphi_2023,ndukuba_resilienciainstitucional_2025}.

No contexto público, a priorização precisa ser simultaneamente técnica, orçamentária e auditável. Registros de ordens de serviço e séries de custos podem apoiar diagnóstico, previsão e planejamento, mas sua utilização depende de qualidade cadastral, classificação consistente e vínculo com ativos, ambientes e intervenções \parencite{martins_custosmanutencaoses_2024,morais_eficienciamanutencao_2023}. A aderência à gestão pública também exige justificativas transparentes para os pesos atribuídos, rastreabilidade das fontes, tratamento explícito de riscos e possibilidade de revisão das decisões. Desse modo, o apoio multicritério é mais útil quando organiza escolhas verificáveis do que quando apenas produz uma classificação final.

A relação entre sustentabilidade e decisão aparece em dois níveis. No primeiro, as dimensões ambiental, econômica, social, técnica-operacional e institucional ampliam o conjunto de efeitos considerados. No segundo, métodos de apoio à decisão estruturam comparações entre alternativas, ativos ou intervenções. A integração entre esses níveis ainda é desigual: são frequentes as menções a sustentabilidade e estrutura decisória, mas menos frequentes as aplicações que demonstram simultaneamente definição de indicadores, ponderação, validação e uso em contexto institucional real. A presença explícita dos ODS em apenas um registro e a ausência da sigla ESG nos campos auditados reforçam a distância entre a abrangência substantiva das dimensões e o uso desses referenciais como linguagem operacional \parencite{masmoudi_ahptopsis_2026,ramantswana_esgfm_2026}.

As coocorrências da matriz analítica mostram quais dimensões, critérios e instrumentos foram codificados simultaneamente nos mesmos registros, mas não estimam associação estatística, causalidade, importância ou peso normativo. Sua contribuição é organizar um espaço de decisão informado pela literatura. Para uso institucional, será necessário converter cada grupo em indicadores com fonte, unidade, periodicidade, direção de preferência e regra de normalização, além de definir como usuários, equipes técnicas e gestores participarão da ponderação.

A força das conclusões é condicionada pela heterogeneidade dos desenhos e pelo nível predominantemente documental da revisão. Estudos conceituais, revisões, aplicações quantitativas, estudos de caso e modelos digitais respondem a perguntas distintas e não oferecem o mesmo tipo de evidência. Como não houve avaliação metodológica baseada no texto completo dos 104 estudos, as frequências não sustentam recomendação sobre superioridade de método, critério ou tecnologia. O uso pontual de texto completo qualificou exemplos específicos, mas não modifica esse limite geral.

A agenda de pesquisa envolve operacionalizar a matriz com indicadores verificáveis; testar métodos de ponderação e agregação com análise de sensibilidade; validar a estrutura em portfólios públicos universitários; e avaliar os efeitos das decisões sobre custos, riscos, desempenho, consumo de recursos e experiência dos usuários. Estudos futuros também devem comparar resultados entre tipologias, explicitar a diferença entre método mencionado e aplicado e documentar a participação dos atores institucionais. Esse percurso permite avançar de uma síntese documental para uma governança preditiva e auditável da manutenção, sem atribuir à matriz um grau de validação que ainda não foi demonstrado.
